%----------------------------------------------------------------------------------------
%	PACKAGES AND OTHER DOCUMENT CONFIGURATIONS
%----------------------------------------------------------------------------------------

\documentclass[10pt,stdletter,dateno,sigleft]{newlfm} % Extra options: 'sigleft' for a left-aligned signature, 'stdletternofrom' to remove the from address, 'letterpaper' for US letter paper - consult the newlfm class manual for more options

%\usepackage{charter} % Use the Charter font for the document text
\usepackage{graphics}

\usepackage{standalone}

\usepackage{hyperref}
\usepackage[american]{babel}
\usepackage{csquotes}
\usepackage[backend=biber,style=authoryear,date=year,dashed=false,doi=false,isbn=false,url=false,arxiv=false]{biblatex}
\renewbibmacro{in:}{} % this should go immediately below calling biblatex 
%\DeclareLanguageMapping{american}{american-apa}
\addbibresource{/Users/hectorbahamonde/Bibliografia_PoliSci/library.bib} 
\hypersetup{
    bookmarks=true,         % show bookmarks bar?
    unicode=false,          % non-Latin characters in Acrobat's bookmarks
    pdftoolbar=true,        % show Acrobat's toolbar?
    pdfmenubar=true,        % show Acrobat's menu?
    pdffitwindow=true,     % window fit to page when opened
    pdfstartview={FitH},    % fits the width of the page to the window
    pdftitle={My title},    % title
    pdfauthor={Author},     % author
    pdfsubject={Subject},   % subject of the document
    pdfcreator={Creator},   % creator of the document
    pdfproducer={Producer}, % producer of the document
    pdfkeywords={keyword1} {key2} {key3}, % list of keywords
    pdfnewwindow=true,      % links in new window
    colorlinks=true,       % false: boxed links; true: colored links
    linkcolor=blue,          % color of internal links (change box color with linkbordercolor)
    citecolor=blue,        % color of links to bibliography
    filecolor=blue,      % color of file links
    urlcolor=blue           % color of external links
}

%% with this patch, the newlfm should work (BEGIN)
\usepackage{etoolbox}
\makeatletter
\patchcmd{\@zfancyhead}{\fancy@reset}{\f@nch@reset}{}{}
\patchcmd{\@set@em@up}{\f@ncyolh}{\f@nch@olh}{}{}
\patchcmd{\@set@em@up}{\f@ncyolh}{\f@nch@olh}{}{}
\patchcmd{\@set@em@up}{\f@ncyorh}{\f@nch@orh}{}{}
\makeatother
%% with this patch, the newlfm should work (END)


\newsavebox{\Luiuc}\sbox{\Luiuc}{\parbox[b]{1in}{\vspace{0.35in}
\includegraphics[width=1.2\linewidth]{/Users/hectorbahamonde/Job_Market/logo.png}}} % Company/institution logo at the top left of the page
\makeletterhead{Uiuc}{\Lheader{\usebox{\Luiuc}}}

\newlfmP{sigsize=50pt} % Slightly decrease the height of the signature field
\newlfmP{addrfromphone} % Print a phone number under the sender's address
\newlfmP{addrfromemail} % Print an email address under the sender's address
\PhrPhone{\texttt{p}} % Customize the "Telephone" text
\PhrEmail{\texttt{e}} % Customize the "E-mail" text

\lthUiuc % Print the company/institution logo

% email input begin
\newread\fid
\newcommand{\readfile}[1]% #1 = filename
{\bgroup
  \endlinechar=-1
  \openin\fid=#1
  \read\fid to\filetext
  \loop\ifx\empty\filetext\relax% skip over comments
    \read\fid to\filetext
  \repeat
  \closein\fid
  \global\let\filetext=\filetext
\egroup}
%\readfile{/Users/hectorbahamonde/Bibliografia_PoliSci/email.txt}
% email input end

\usepackage{color} % for colors
\usepackage{hyperref} % for links

\hypersetup{
  colorlinks = true,
  urlcolor = blue,
  pdfpagemode = UseNone
}

%----------------------------------------------------------------------------------------
%	YOUR NAME AND CONTACT INFORMATION
%----------------------------------------------------------------------------------------

\namefrom{\vspace{-3cm}Hector Bahamonde, PhD} % Name

\addrfrom{
{\vspace{-2cm}}\\ % Date
[12pt]
{\bf Hector Bahamonde, PhD}\\
{\color{blue}{\bf Senior Researcher}}\\
\emph{{\color{blue}INVEST Flagship Centre}}\\
{\color{blue}University of Turku}\\
Turku, Finland
}

\phonefrom{+358403621438} % Phone number

\emailfrom{\href{mailto:hibano@utu.fi}{hibano@utu.fi} \\ 
\texttt{w}: \href{http://www.hectorbahamonde.com}{www.HectorBahamonde.com}\\
\today}% Link to last version


%----------------------------------------------------------------------------------------
%	ADDRESSEE AND GREETING/CLOSING
%----------------------------------------------------------------------------------------

\greetto{\vspace{-1cm}Dear Rector of the University of Turku,} % Greeting text
%\closeline{\vspace{-1cm}\includegraphics[width=3cm]{/Users/hectorbahamonde/Administracion/signature.pdf} \hspace{-4cm}} % Closing text




%----------------------------------------------------------------------------------------

\begin{document}
\begin{newlfm}

%----------------------------------------------------------------------------------------
%	LETTER CONTENT
%----------------------------------------------------------------------------------------
 %introduction
\vspace{-0.3cm}I am writing to apply for the title of Docent at the Faculty of Social Sciences of the University of Turku. Currently, I am a \href{https://www.utu.fi/en/people/hector-bahamonde-norambuena}{Senior Researcher} at the INVEST Flagship Centre at the University of Turku. I also serve as the Research Director of the \href{https://www.utu.fi/en/investhub}{INVEST-Hub}, an interdisciplinary center about experimental and quantitative social research. I obtained my B.A. in Political Science at Catholic University of Chile. After receiving my PhD in Political Science from Rutgers University-New Brunswick, NJ, U.S. in May 2017, I spent one year as a post-doctoral fellow at Tulane University-New Orleans, LA, U.S. Before coming to the University of Turku, I was an Assistant Professor (tenure-track) at O$'$Higgins University, in Chile for over two years. My long-term plan is to settle in Turku, and build an academic career at the University of Turku, hopefully, securing a tenure-track position in my field. Thus, the title of Docent, to me, is a stepping stone in my long-term development plan. In this statement I will summarize my accomplishments and future research and teaching plans.

{\bf Research and scientific activities}. Within Political Science, my primary subfield is the political economy of inequality, democratic development and clientelism. In simple, I am interested in how economic conditions---for example, inequality, economic growth, poverty, economic losses and gains, among others---affect democracy itself and people's attitudes toward democracy. Research methods are an important part of my research. Every piece I have published pays considerable attention to both the methodological \emph{and} substantive aspects of the question at hand. Importantly, since I defended my dissertation in 2017, and in my humble opinion, I have been able to---at least---produce the equivalent amount of work corresponding to a doctoral dissertation. The following paragraphs briefly describe (a selection of) my peer-reviewed publications.


In \href{https://doi.org/10.1177/1866802X1801000202}{\textcite{Bahamonde2018}}, and in the context of clientelism in Brazil, I asked whether political parties target individuals or groups. This question is rather important because the incentives, costs and risks of either strategy are quite different. For example, individual defection costs significantly less than group defection. Yet, the literature had neglected that question. Using causal inference tools (matching), the paper finds that parties target individuals when they are identifiable, and groups when parties need to rely on the spillover effects of clientelism. The argument, in simple, exploits the economic concept of marginal costs, but applied to the returns of vote buying. 

Next, in \href{https://doi.org/10.1016/j.ejpoleco.2021.102048}{\textcite{Bahamonde2021a}} we asked about the relationship between democracy, state capacity and economic inequality. If we follow the literature in political economy, we should see that most arguments explain that there should a negative association: the median voter should vote for candidates than maximize their welfare---embodied in redistributive policies. However, using time-series and cross-sectional methods, we find that higher levels of democracy and state capacity \emph{increase} inequality overtime. We exploited a large dataset comprising all countries since the 1970s. Importantly, the causal mechanism operates via the capacity a state of enforcing contracts and protecting property rights, increasing the returns of the banking sector, thus, increasing inequality overtime. This piece is rather important also because it was \href{https://www.nytimes.com/2021/08/04/opinion/biden-eviction-covid-democrats.html}{quoted} in the New York Times by one of the most influential columnists of this important newspaper.

After that, in \href{https://link.springer.com/article/10.1057/s41269-020-00174-4}{\textcite{Bahamonde2020a}}, I published the results of a novel list experiment fielded in the United States. Unlike other research questions that ask factual questions, in this piece I asked whether voters \emph{would} vote for a candidate in exchange for money. Thought clientelism is low in the U.S., that does not necessarily mean that voters \emph{would not} sell their vote to a political party. It is important to mention that list experiments are excellent tools to study questions loaded with social desirability bias (such as selling one's vote). I find in the paper that 25\% of voters---it was a representative sample---would vote for a candidate in exchange for money. The paper teaches an important lesson: while the levels of clientelism are comparatively low in the U.S., scoring ``high'' in most democracy indexes does not necessarily imply that American voters will not sell their right to vote freely. Coincidentally, the experimental fielding took place during the political campaign that gave Trump the presidency of the U.S. As I explain below, I have a work-in-progress manuscript where I exploit this unique election, which I consider a critical juncture.  


The same year we published in \href{https://doi.org/10.1016/j.electstud.2022.102497}{\textcite{Bahamonde2022b}} the results of a laboratory experiment. In particular, we asked when parties buy more votes, whether it is when they are winning or losing the election. From a substantive point of view, this is an important question. The literature argues something that makes sense: parties should buy more votes when they are losing an election. If we assume that parties want to remain in power, spending on vote buying when losing the election seems an optimal strategy. However, we discovered that parties engage in clientelism more when they are \emph{winning} the election because parties are more concerned with preventing a decline rather than increasing gains. From a theoretical standpoint, we theorized from a prospect theory approach. Empirically, we find that the more electoral assets parties acquire during their campaigns (i.e., supporters/voters), the more there is to lose, making them to invest more in clientelism in favorable electoral contexts. 

Finally, I would like to take this opportunity to talk about a paper that received a \emph{Revise and Resubmit} in Political Psychology. The paper is in its second round of revisions, and after seeing the reviews, we are quite confident that that paper will be published. In the context of a natural experiment in Finland, we discovered in \textcite{Bahamonde:2023} that candidates that look like working class have less chances of winning an election, particularly female candidates. This piece refines a very old empirical regularity: attractive candidates win elections more often. However, the literature had been silent about a different yet important question: \emph{What happens when a candidate looks poor/wealthy? Does that affect the candidate's probability of being elected?} We discovered that there is an important the electoral penalty for looking poor. Female candidates are the ones that get more penalized by their ``working class'' looks. We are quite certain this paper will give lots to talk about because Finland is usually characterized as an equal society, consequently, making it a hard case for unequal outcomes. We pretend to go public with these findings, particularly now that there are elections in Finland.

As it can be seen, my understanding of the political economy of voting behavior, inequality and clientelism, departs from neo-classical approaches. In fact, my starting point is that agents (e.g., voters, parties and states) have a hard time assessing what is ``best'' or ``optimal'' for them, and rather take cognitive shortcuts that help them to deal with their respective realities. In addition to that, and from a quantitative and qualitative standpoints, I have been quite active publishing in some of the top journals in the discipline. In fact, I will take the liberty to argue as well that throughout this time, we had three children with my wife. My oldest son was just a few months old when he sat in my dissertation defense. Then, as it can be seen in my CV, several movings happened (within the U.S. and abroad as well). Also, then the COVID pandemic hit, slowing me down even more, particularly with the collection of the the lab experiment paper. However, I'm a very determined person and managed to accomplish this list of publications, of which I am very proud of. 

{\bf Next steps} in my research plan include to keep working on several projects. First, I have a novel dataset that I collected while working on \textcite{Bahamonde2020a}. These data are a conjoint experiment where I find that Democrats in the United States stick slightly less to democratic principles compared to Republicans. These data are a unique piece of evidence as it was (coincidentally) collected during the campaign that gave Trump the Presidency of the U.S. My working hypothesis is that Democrats have a looser adherence to authoritative institutions (such as Congress and the Presidency), making them less strict about their democratic commitments. More importantly, I am currently working on my application to get funding from the \emph{Academy of Finland}. My project seeks to investigate---throughout several survey and lab experiments---how inequality shapes attitudes toward democracy. This question is particularly relevant in the current context of democratic erosion in many of the countries we once considered ``strongly democratic.''

Moving forward, my research philosophy is to advance some of the most classic questions in comparative politics. Thus, questions about democratic development and democratic undermining are the ones that drive my intellectual curiosity. In particular, my future research intends to investigate the conditions that lead people to support anti-democratic candidates and policies; this is within the scope of the literature on democratic backsliding. 

In addition to that, I always approach these questions with a strong methodological component. I'm not geographically constrained, thus, my papers study different and diverse countries such as Brazil, Finland, and the United States. However, one aspect that is present throughout all my research is the adherence to the highest ethical standards. Since often times I employ experimental methods, it is paramount to me to always be transparent with the participants I recruit for my studies. Also, I am true believer in scientific reproducibility: every paper I have published has its own GitHub repository where not only the data but also the code I used is publicly available. 

Finally, another aspect that takes an important part in my academic activities, is community building. It is not unusual that I organize panels or entire sections for different conferences. While the details are in my CV, I would like to name an example: I (along with my colleague from Abo Akademy) put together an entire track on ``democratic backsliding'' for the last Finnish Political Science Association. In addition to that, I usually review for some of the most prestigious outlets in the discipline. I am also an \href{https://www.springer.com/journal/43545/editors}{Associate Editor of Social Science}, a Springer Nature journal.

{\bf Teaching}. I have considerable experience in teaching. First, I obtained a fellowship from Rutgers University that covered my entire doctoral studies (plus stipends), for which in return I had to be a Teaching Assistant for three years. I quickly became a really good teaching assistant. My duties were actual teaching (lecturing), grading and test designing. Then, I was the teaching assistant for the Introduction to Quantitative Methods, a course all PhD students had to take in my program. Later on, when I was a post-doctoral fellow in Tulane University, I designed and taught my first course. It was ``Introduction to Comparative Politics.'' It was such a great experience, and the students evaluated me very positively (teaching evaluations are available upon request). I feel that I really consolidated my teaching skills when I moved back to Chile, where I was an assistant professor for a little more than two years. On top of my research duties, I had to teach between two and three courses per semester. Naturally I took over the epistemology, research design and statistics courses. But also, I taught the entire comparative politics sequence. I have had similar substantive/methodological teaching responsibilities at the University of Turku (but in lesser quantity). I could say that teaching Masters students for the first time taught me a very important lesson: students come with different substantive (and methodological) backgrounds and interests. Thus, I quickly adapted to those, and assigned papers and gave examples accordingly. For example, some students were already working on their respective research projects, thus, I always tried to give examples that resonated better with those projects.

% Pedagogical approach and training
I love teaching. My approach is to challenge students, always trying to stimulate their intellectual curiosity. For example, rather than presenting the assigned readings as the cannon in the subfield, my approach is to make my students to be merciless in their critical approach to those readings. Like I always tell them, I do not want my students to be obedient; more on the contrary, I want them to learn to be critical thinkers, always ready to give constructive criticism. \emph{Is the case/methodology appropriate to answer the question at hand? Are the argument and the empirical evidence sound? What would \underline{you} do differently?} Those are some of the questions I tell my students to focus on when they are approaching the readings. That is why I usually assign reaction papers that are due \emph{before} my lecture. Office hours are really important to me, and my doors are always open to answer question. I contribute to a variety of needs: mentoring needs, course material needs, substantive and methodological needs. {\bf I like to spend quality time with my students---I certainly valued my former instructors/professors when they were around and ready to solve doubts}. 

It is my plan to have more teaching responsibilities. For example, I would like to teach substantive and methodological courses, particularly, for Masters and PhD students. Supervising theses and sitting in dissertation committees are also part of the experience I would like to have at the moment. I feel ready. I think I can contribute a big deal, especially, to the comparative politics and methods area. I would feel very comfortable teaching the Comparative Politics and the entire methods tracks. I would be very happy to take care of the departmental needs as well. I am a very flexible person. Some of the courses I would like to teach are (the list is flexible and there is considerable room for other options as well, depending on the Department's needs): Comparative Politics I and II, Political Economy, Research Design and Epistemology in the Social Sciences, Quantitative Methods I (linear regression), Quantitative Methods II (generalized linear models via maximum likelihood), Experimental Methods (natural, laboratory and survey experiments). It is very important to clarify that my approach to teach all methods courses is very practical. We do talk about some statistical theory, but the focus is on empirical applications. We always use data and software to make these matters more concrete and approachable.

{\bf Administrative tasks and other activities}. I feel comfortable contributing to the administration of the institutions I have worked for. For example, I was appointed Program Director of the Public Administration program when I was in Chile. My job consisted of making sure that all teaching needs of the students of the program were covered. In practice that meant, to hire instructors, to modify the study plan, to coordinate extra-curricular activities, to mentor students, to give public speeches, and to represent the interests of the program with the Rector of the University. Later on, at the University of Turku, I was appointed as the Research Director of the INVESThub, an interdisciplinary instance of INVEST focused on experimental and quasi-experimental research. My main task consists of building a solid community of researchers working on diverse substantive topics, but all of them using experimental methods. As the Director, I am in charge of the Seminars series, where different researchers present their work. We also organize conference panels and currently are evaluating our teaching contribution to the INVEST centre. Importantly, I make sure our technical needs are covered. For example, we recently hired an IT person in charge of the servers (used for several experiments).

As a Docent at the University of Turku, I would be very happy to take over of as many administrative tasks as needed, such as, teaching coordination and the speakers series of the Department. I have performed these tasks before, have done it well, and would be very happy to do it here at the University of Turku.

Finally, I think our role in academia does not end within the University walls. That is why I always try to contribute as much as possible to society. While the details are in my CV, I am constantly writing op-eds. Recently, I gave an in-depth \href{https://twitter.com/CNNChile/status/1527122195780186112?s=20&t=-tO34V0-IHexWmRUePwHJ/}{interview to CNN-Chile} about Finland and its relationship with NATO. 

{\bf To summarize}, I would


% ending
{More information, \href{http://www.hectorbahamonde.com/teaching/}{\texttt{syllabi}}, my \href{http://github.com/hbahamonde/Job_Market/raw/master/Bahamonde_Research_Statement.pdf}{\texttt{research}}, \href{http://github.com/hbahamonde/Job_Market/raw/master/Bahamonde_Teaching_Statement.pdf}{\texttt{teaching}} and \href{http://github.com/hbahamonde/Job_Market/raw/master/Bahamonde_Diversity_Statement.pdf}{\texttt{diversity}} statements, as well as other \href{http://www.hectorbahamonde.com/research/}{\texttt{papers}} are available on my website: \href{http://www.hectorbahamonde.com}{\texttt{www.HectorBahamonde.com}}. Thank you for considering my application. I look forward to hearing from you.\unskip}

% References
\newpage
\printbibliography


%----------------------------------------------------------------------------------------

\end{newlfm}
\end{document}

The academic portfolio is structured as follows:
#. Basic information
#- Personal and contact information
#- Education and degrees
#- Other relevant background information
#2. Research and academic activities
#- Experience in research and other scientific activities
# - Research philosophy and ethics, primary fields of research
#- Most notable publications
- Research evaluations and awards
#- Activity in the academic community
- Vision and personal development
- Other scientific qualifications
3. Teaching and supervision
# - Experience of undergraduate and doctoral teaching and supervision
#- Pedagogical approach and training
- Development of teaching skills and awards for teaching
- Participation in the assessment and development of teaching
- Strengths, development challenges and visions for own teaching
- Other qualifications related to teaching
4. Administrative tasks and other activities
#- Administrative and management duties
- Positions within the field outside academia
- Societal activities and positions of trust
- Publications, presentations and the dissemination of scientific knowledge
through other means
- Other significant qualifications and tasks
